\section{Navigation Architecture and State Estimation}
\label{sec:navigation}

The navigation filter, \texttt{kalman.py}, is a multiplicative
error-state \ac{MEKF}: inertial prediction on the \ac{IMU} stream,
corrected by the \ac{GNSS}, \ac{LIDAR} and star-tracker measurements of
Section~\ref{sec:sensors}. The estimated state is referenced to the \ac{IMU}
station --- its position and velocity in the external frame --- so the
lever arms of the other sensors are constant Body vectors and no state
re-referencing occurs when the center of gravity migrates during burn.
The nominal 24-state configuration is
$\mathbf{x} = [\mathbf{p}, \mathbf{v}, \mathbf{b}_a,\,
\mathbf{q}, \mathbf{b}_g,\, \mathbf{s}_f,\, \mathbf{m}]$:
\ac{IMU}-station position and velocity, accelerometer bias, the attitude
quaternion from Body to \ac{ECEF} (Hamilton, scalar-first), gyro bias,
per-axis gyro scale factor and the six off-diagonal gyro misalignment
terms (Table~\ref{tab:nav-state}).

\begin{table}[htbp]
\centering
\begin{tabular}{lll}
\toprule
Block & Content & Initial $1\sigma$ \\
\midrule
$\mathbf{p}$, $\mathbf{v}$ & \ac{IMU}-station position, velocity [\textrm{m}, \textrm{m/s}] & $2\,\mathrm{m}$, $0.05\,\mathrm{m/s}$ \\
$\mathbf{b}_a$ & accelerometer bias, Body [$\mathrm{m/s^2}$] & $0.1$ \\
$\mathbf{q}$ & attitude (multiplicative error $\delta\boldsymbol{\theta}$, Body [rad]) & $0.01\,\mathrm{rad}$ \\
$\mathbf{b}_g$ & gyro bias, Body [$\mathrm{rad/s}$] & $10^{-3}$ \\
$\mathbf{s}_f$, $\mathbf{m}$ & gyro scale factor, misalignment & $10^{-3}$ \\
\bottomrule
\end{tabular}
\caption{Estimated-state blocks and initial standard deviations. The
initial position error is a seeded $[1, -1, 0.5]\,\mathrm{m}$; the attitude
starts at truth with a $0.01\,\mathrm{rad}$ prior unless declared
otherwise.}
\label{tab:nav-state}
\end{table}

\subsection{Prediction}

Each $100\,\mathrm{Hz}$ \ac{IMU} sample advances the strapdown state. The
gyro measurement is corrected with the current calibration estimate,
$\hat{\boldsymbol{\omega}} = \hat{\mathbf{G}}^{-1}
(\boldsymbol{\omega}_{\mathrm{meas}} - \hat{\mathbf{b}}_g)$, and the
attitude is propagated by its quaternion transport equation with the
trapezoidal mean rate over the interval; the mid-interval attitude
rotates the bias-corrected specific force into the external frame
alongside the modelled gravity and rotating-frame accelerations of
Section~\ref{sec:environment}, whose analytical gradients linearize that
model in the filter. The error dynamics couple the multiplicative attitude
error into the velocity error through $-\mathbf{R}\,[\hat{\mathbf{f}}]_\times
\delta\boldsymbol{\theta}$, and the gyro chain drives the attitude error
through the bias, scale-factor and misalignment blocks --- the calibration
is observable only through attitude aiding and sufficient angular motion.
Process noise is built from the sensor \acp{PSD}
($\sigma_a^2/f_{\mathrm{IMU}}$, $\sigma_g^2/f_{\mathrm{IMU}}$) and the
declared bias random walks.

\subsection{Measurement updates}

All corrections use the Joseph-form covariance update and record a
normalized innovation square for monitoring. The \ac{GNSS} update corrects
position and velocity: its $0.1\,\mathrm{s}$ latency is compensated by
extrapolating the measurement to the current epoch with the estimated
acceleration, its lever arm (antenna minus \ac{IMU} station) is removed
with the angular rate, and the measurement Jacobian carries the full
attitude, bias and calibration sensitivity of that arm. The \ac{LIDAR}
update is scalar and nonlinear: the measurement is the analytic
line--sphere slant range of Section~\ref{sec:sensors}, differentiated
analytically at the estimated state, with attitude sensitivity through
both the arm and the boresight direction; invalid geometry skips the
update. The star-tracker update is an absolute attitude correction: the
measurement quaternion is converted from \ac{ECI} to \ac{ECEF} by the
Earth rotation at the exposure epoch, propagated over its latency with
the last estimated body rate, and differenced multiplicatively from the
estimated attitude; its noise covariance is rotated into the Body frame by
the known mounting, and the update is gated by the three-degree-of-freedom
\ac{NIS} at the $99.9\,\%$ level ($16.27$), with rejects counted rather
than silently absorbed.

\subsection{Interfaces and diagnostics}

The filter consumes its own attitude estimate by default
(\texttt{use\_estimated\_attitude}): the prediction rotates the specific
force with the estimated quaternion, and every GNSS, \ac{LIDAR} and
star-tracker correction also corrects that attitude. A declared
diagnostic mode retains the original nine-state filter with the truth
attitude supplied from outside; its results are kept distinct from the
gyro-driven results in the results chapter. A second, deliberately
unaided baseline propagates the attitude from the raw gyro stream alone
(\texttt{GyroAttitude}): starting from the same initial attitude, it
carries uncompensated gyro noise and bias, and its logged error isolates
how much the aiding buys over pure dead reckoning. One filter instance
runs across staging without a reset --- the estimated state is
\ac{IMU}-station referenced and the sensor stations are stage-1
hardware, so the booster return needs no state handover; the separated
upper stage flies an open-loop truth continuation without a navigation
instance.

\subsection{Verification and nominal-run behavior}

The filter is verified in stages: numeric finite-difference checks confirm
the analytic Jacobians of the attitude error with respect to gyro bias,
scale factor and misalignment; a closed-loop scenario (a vehicle
exercised on the rotating Earth with star-tracker aiding) recovers
injected bias, scale-factor and misalignment values; and the cross-module
and mission-level suites (115 tests, all green) exercise the integrated
chain end to end. In the nominal full-mission run the filter operates in
its default gyro-driven configuration: the position error holds
$0.881\,\mathrm{m}$ \ac{RMS} ($0.918\,\mathrm{m}$ in the truth-attitude
reference configuration), the aided attitude error stays at
$0.456^{\circ}$ \ac{RMS} (peak $1.57^{\circ}$, touchdown
$0.012^{\circ}$) while the unaided gyro baseline drifts monotonically to
$1.89^{\circ}$, and the $z$-axis gyro bias is recovered to $1.4\,\%$.
The \ac{GNSS} innovations are statistically consistent (mean \ac{NIS}
$7.58$ for six degrees of freedom) and \ac{LIDAR} delivers $496$ valid
updates. All $381$ valid star-tracker exposures are rejected by the
\ac{NIS} gate: the propagated attitude uncertainty ($1\sigma$ median
$0.118^{\circ}$) understates the actual maneuver-induced error by more
than the factor the $99.9\,\%$ gate tolerates, so the nominal attitude is
gyro-propagated with weak \ac{GNSS} lever-arm aiding only --- the gate
outcome and its causes are treated as a result in the results chapter of
this report.

\paragraph*{Assumptions and limitations}
All noise and bias parameters are the placeholder research assumptions of
Table~\ref{tab:sensors}; the gyro calibration block estimates a unit whose
adopted true misalignment and scale factor are zero, so its output is a
zero-check rather than a calibration demo; \ac{GNSS} latency compensation
assumes constant acceleration across the $0.1\,\mathrm{s}$ window;
attitude prior and initialization are declared values, not surveyed
alignment; and the filter has no protection against unmodeled faults
beyond the star-tracker \ac{NIS} gate --- every update it accepts is
trusted once inside the gate.
